% Pursuing Stacks (cotes 134-2 to 134-8): archive cote 134-2, site page 61 (PDF page 62).
% appendix (letters to Breen). Draft compared: ps1.tex:1439-1461.
% AI-transcribed by Claude Opus 5.5 (claude-opus-5-5, adaptive thinking, effort high), Message Batches API, batch ps1 (msgbatch_01Rv5hCaNTtUt4JHx1yqoHPL), 2026-11-12; not yet checked by eye.
\scanpage{134-2/61}
\mspage{-46 -}

J$_0$ if I remember well, (J$_0$ had as abelian part the abelian part of A\hand{lb}$_{X/k}$\hand{\unsure{)|,}}
the \deleted{ordinary}\inserted{usual} generalised jacobian). It\hand{'}s construction, inspired by the
residual complex, passes by generalised jacobians (in an appropriate cohomological
sense\hand{)} of the \deleted{localised}\inserted{localizations} Spec $\underline{O}_{X,\ x}$ of $X$ at its different point.
\margin{$\llcorner$}N.B. $\underline{H}_0$ $(J_*)$ was the ``generalised Jacobian'' of $X$ , i.e. there existed a
homeomorphism $X \to \underline{H}_0$ , which was universal for homomorphisms of $X$ into \hand{comm.}
locally proalgebraic groups. For $X$ connected, $\underline{H}_0$ is an extension of $\mathbb{Z}$
by an appropriate proalgebraic group.
It is possible that\hand{,} at first, I restricted to the case of $X$ smooth .
The cohomology role of this complex was that of a complex of \textit{homology}
\margin{formule sur une ligne}(*) $H^i(X, G_X) \simeq \mathbb{E}xt^i(J_{*X/k}, G)$ but for which coefficients\hand{\unsure{$G$?}} I believe
I took arbitrary commutative algebraic groups $G$ but worked with the
Zariski topology (malédiction !)\hand{.} \over{e}{E}ven in the case of discrete $G$ , I considered
the Zariskian $H^i$ , this gives slightly stupid cohomology groups, evidently.
I realised that one should work ultimately in étale cohomology, and that
the construction of the \hand{(}$J$\hand{)}$_{\unsure{i}X/k}$ will evidently be modified\inserted{accordingly}. As for the
significance of the $\mathbb{E}xt^i$ (hypercohomology), at a moment where Serre had
developed the formalism for proalgebraic groups, one was not too fearful of
taking it in the category of such objects - \deleted{the case of}\inserted{and \deleted{\illegible{1}} in the sense of a} ``derived category''
which at that moment had never yet been \inserted{explicitly} defined \deleted{or}\inserted{and} studied\hand{.}(We have, after
all, somewhat progressed since those days!). I have the impression, in
view of these antique cogitations, heuristic as they were, that it should
now be possible to develop at present such a theory of $J_{*X/k}$ , in cohomology
\margin{\illegible{1}}f.p.p.f., giving a formula (*) without limitation on the degree $i$ of the
cohomology. (N.B. But $J_*$ evidently no longer stops in dim $X = n$ but
in dim $2n$ . It is nevertheless possible that the components $J_i$ might be
of dim 0 for $i > n$ ) .

I believe that the construction of the $J_*$ does not commute with base
change, but merely does so in the derived category sense.

* *
*

\margin{laisser espace}
\margin{61}
